\scanpage{59}
% archive foliation: 58
% lines 1--2
\[
f\longmapsto\left(
\left(\underbrace{\unsure{f_0(u_x)|f(u_x)}}_{\ell_x}\right)_{x\in C_0},
\left(f_{a_i,a_i}:\underbrace{\operatorname{Aut}(a_i)\longrightarrow
\unsure{\operatorname{Aut}}(f_0(a_i))}_{\varphi_i}\right)\right).
\]
% lines 3--4
c) Pour que $f$ soit fidèle, il f. et suffit que $\forall i\in I$, $\varphi_i$ soit injectif.

% lines 5--8
Pour que $f$ soit pleinement fidèle, il f. et suffit que ;
$1^\circ)$ pour $i\ne j$, on ait
$\operatorname{Hom}(f_0(a_i),f_0(a_j))=\varnothing$

% lines 9--10
$2^\circ)$ pour $\unsure{\forall i|\forall}\ x,y\in C(i)_0$,
\[
\unsure{\operatorname{Isom}}(f_0(x),f_0(y))\overset{\sim}{\longrightarrow}
\operatorname{Hom}(f_0(x),f_0(y))
\]
% line 11
$3^\circ)$ $\forall i\in I$, $\varphi_i$ est un iso.

% lines 12--16
\deleted{\unsure{d)} Si $C$}

Pour que $\unsure{f}$ soit une équivalence,
il f. et s. qu'on ait de plus

$4^\circ)$ Tout $x'\in C'_0$ est isomorphe à un des $f_0(a_i)$

% lines 17--21
Pour que $f$ soit un isom.,
il \unsure{f.} et s. qu'on ait $1^\circ$, $2^\circ$, $3^\circ$,
et de plus

$4'^\circ)$ $f_0$ bijective.
